\chapter{State of the Art}
\label{cap:stato_arte}

This chapter reviews vision-based methods for agricultural row
navigation, focusing on the conversion of image observations into a
local guidance reference. As surveyed by Bai et al.~\cite{bai2023vision},
vision sensors provide rich appearance and geometry cues for autonomous
field guidance at low hardware cost. The review covers visual perception,
geometric extraction, and the integration of these stages with robot control.
Dataset preparation and the implementation used in this thesis are
described in Chapters~\ref{cap:segmentation}
and~\ref{cap:geometric_reference}.

\section{Classical image-processing methods}
\label{sec:classical_methods}

Classical pipelines combine image filtering, thresholding, and
application-specific geometric rules. Mendez et al. convert vineyard
images to grayscale, apply Gaussian filtering and Otsu thresholding,
select the largest blob, and use column pixel counts to estimate the row
center \cite[Section~3.4]{mendez2023vineyards}. Such methods can be
effective in the conditions for which their parameters were selected,
but changes in illumination, weeds, and crop appearance can require
retuning \cite{bai2023vision}.

\section{Learning-based segmentation for agricultural navigation}
\label{sec:learning_based_segmentation}

\subsection{Semantic segmentation approaches}
\label{subsec:semantic_segmentation}

Semantic segmentation assigns a class to each pixel, but the label
definition determines what the geometric stage can recover. De Silva
et al.\ train U-Net to predict narrow crop-row masks in sugar-beet
images. Their labels represent row structure rather than individual plant
pixels or traversable soil~\cite[Sections~3.3 and~4.1]{silva2023rowdetection}.
Navone et al. instead predict a binary vegetation mask from RGB images
using a MobileNetV3 backbone and an LR-ASPP head, then combine that
output with depth for reference extraction
\cite{navone2023ecmr,navone2025gpsfree}.

Like other path-oriented approaches in the literature, this work segments
the traversable soil corridor between crop rows rather than the plants
themselves. The segmentation model provides a candidate region of drivable
ground, and the downstream geometric stage then extracts the lateral
boundaries and computes the steering reference.

\subsection{Training data and generalization}
\label{subsec:training_data}

Pixel-level annotation is time-consuming, and agricultural datasets
must represent changes in crop, growth stage, and lighting. De Silva et
al. include 11 field variations in their dataset
\cite[Section~3]{silva2023rowdetection}. Navone et al. initialize their
segmentation model with ImageNet weights and train using synthetic
AgriSeg images supplemented by 100 real crop images
\cite[Section~4.1, Table~1]{navone2025gpsfree}. Pretraining and the
choice of training images are separate design choices.

\section{Semantic segmentation architectures}
\label{sec:segmentation_architectures}

\subsection{FCN, DeepLabV3, and mobile architectures}
\label{subsec:architectures_overview}

Fully Convolutional Networks (FCN) replace fully connected layers with convolutions to produce dense pixel-level predictions~\cite{long2015fcn}. DeepLabV3 incorporates atrous convolutions and atrous spatial pyramid pooling to capture multi-scale context without decreasing feature resolution~\cite{chen2017deeplab}. For embedded applications where compute resources are constrained, Howard et al.~\cite{howard2019mobilenetv3} introduced MobileNetV3 together with the Lite Reduced Atrous Spatial Pyramid Pooling (LR-ASPP) head, which reduces latency by using depthwise separable convolutions.

The choice of architecture affects both segmentation accuracy and inference speed. Deeper backbones or larger receptive fields can yield higher IoU values at the cost of additional compute. For robotic systems operating in closed loop, inference latency on edge hardware dictates the achievable control frequency.

\subsection{Transfer learning and pretraining}
\label{subsec:transfer_learning}

Transfer learning initializes network weights from models pretrained on large visual corpora such as ImageNet. The pretrained backbone is a generic feature extractor, while the task-specific segmentation head is adapted to produce class scores matching the target dataset. In agricultural vision, pretraining reduces training time and mitigates overfitting when domain-specific annotated datasets are limited in size. For instance, Li et al.~\cite[Section~2.2.2]{li2022multiroi} demonstrate that transferring weights learned on corn fields to other crops (such as tomatoes, cucumbers, and wheat) preserves segmentation accuracy with small training sets. Similarly, Navone et al.~\cite[Section~4.1]{navone2025gpsfree} initialize their segmentation backbone with ImageNet weights prior to fine-tuning on tree-row imagery.

The extent of pretraining differs across implementations. Some setups retain pretrained segmentation heads and replace only the final classification layer, whereas others initialize only the backbone and train the segmentation head from scratch.

\section{Dataset preparation and model-assisted annotation}
\label{sec:dataset_annotation}

\subsection{Manual annotation challenges}
\label{subsec:manual_annotation}

Manual pixel-level annotation requires annotators to label every pixel in training images. For agricultural datasets spanning multiple crops, growth stages, and lighting conditions, this process is time-consuming and expensive. Annotation consistency depends on annotator interpretation of ambiguous boundaries, particularly where plants overlap soil or covering material.

\subsection{SAM and foundation models for segmentation}
\label{subsec:sam_annotation}

Foundation models for segmentation have established model-assisted annotation as an effective alternative to purely manual labeling. The Segment Anything project built its dataset with a model-in-the-loop data engine in three stages (assisted-manual, semi-automatic, and fully automatic), obtaining more than one billion masks on 11 million images~\cite{kirillov2023sam}. While early foundation models required spatial prompts such as bounding boxes or foreground/background clicks for every target instance, Segment Anything Model~3 (SAM~3) introduces open-vocabulary concept segmentation driven directly by textual descriptions and visual exemplars~\cite{carion2026sam3}.

In precision agriculture, model-assisted annotation with foundation models has been explored primarily for satellite-based field boundary delineation and canopy phenotyping. However, its systematic application to agricultural ground navigation remains largely unexplored. In this thesis, SAM~3 is employed offline as a zero-shot concept annotator, and its masks are used without manual correction. On lettuce, networks trained on SAM~3 labels with different class sets (path, plants, plastic mulch) are compared against manual ground truth; binary path labels generated for the complete lettuce, tomato, and cucumber collection train the network deployed on the robot. This offline annotation pipeline decouples the heavy computational burden of foundation models from the real-time inference requirements of the onboard robot controller.

\section{Geometric extraction methods}
\label{sec:geometric_extraction}

The geometric stage reduces a mask to a quantity the controller can use.
Depending on the representation and method, the output may be a
horizontal image position, a straight line, a crop-edge estimate, or a
curve. These outputs describe different tasks and coordinate systems.

\subsection{Line fitting and crop-edge detection}
\label{subsec:line_fitting}

Chen et al. estimate a path line for greenhouse cucumber rows from
navigation points selected in plant--soil segmentation. Their
prediction-point Hough transform uses a least-squares estimate to
restrict the line search and is compared with traditional Hough and
least-squares fitting \cite[Sections~2--3]{chen2021navigation}.

Kneip et al. detect crop--ground transitions in scanlines of a stereo
elevation map, then fit a linear crop-edge model using a Huber
M-estimator and iteratively reweighted least squares. Their output is a
crop boundary in ground-plane coordinates, not an image-space estimate
of the inter-row path \cite[Sections~6--7]{kneip2020cropedge}.

\subsection{Column-wise histogram methods}
\label{subsec:histogram_methods}

Histogram methods aggregate mask pixels by image column. Mendez et al.
estimate the row center from the five columns with the largest pixel
counts in a thresholded background region
\cite[Section~3.4.3]{mendez2023vineyards}. Navone et al. use a different
input: a temporally combined vegetation mask refined with depth. SegMin
smooths its column histogram and selects the minimum; SegMinD weights
the vegetation mask by normalized inverse depth before calculating the
histogram \cite[Section~3]{navone2025gpsfree}. In both cases the output
is a horizontal position, not a full path line.

\subsection{Horizontal strips and adaptive multi-ROI methods}
\label{subsec:strip_methods}

Li et al. combine a Faster-U-net inter-row-path mask with an adaptive
multi-ROI procedure. Horizontal regions supply ordered local points,
which are used as control points for a multisegment cubic B-spline.
Their experiments use 40-pixel search regions and discuss the trade-off
between the number of points and fitting behavior
\cite[Section~2.2.3, Figure~9]{li2022multiroi}.

This is the closest geometric comparison for the present work because
both pipelines select local evidence from horizontal image regions.
The representation and fitted model differ: Li et al. use an
inter-row-path mask and a spline, whereas this thesis estimates a local
straight reference from the soil corridor between parallel crop rows.

\subsection{Triangle Scan Method}
\label{subsec:triangle_scan}

De Silva et al. propose the Triangle Scan Method (TSM) for crop-row
masks. It first scans an upper horizontal band to find an anchor column,
using a preset anchor if the scan fails a threshold. Candidate lines
from that anchor to positions along the base of a triangular region are
then evaluated; the line with the greatest mask support defines the
central crop row. Its two endpoints provide position and angular errors
for visual servoing \cite[Section~4.2, Algorithm~1]{silva2023rowdetection}.

TSM is not the same as extracting one point from each of several
horizontal strips: the upper scan establishes an anchor for a bounded
line search. It also operates on crop-row masks, so its output is not
directly equivalent to a reference extracted from a soil mask.

\subsection{Bilateral boundary tracking versus centroid reduction}
\label{subsec:bilateral_vs_centroids}

In agricultural vision pipelines that convert binary segmentations into geometric references, existing methods fall into distinct categories, each characterized by specific operational trade-offs:

\begin{enumerate}
    \item \textbf{Crop-canopy feature extraction with RANSAC:} A classical paradigm detects feature points along crop canopies or plant stalks, partitions them into left and right row clusters, and fits linear models using Random Sample Consensus (RANSAC) combined with ordinary least squares~\cite{bai2023vision}. Recent monocular vision implementations extract crop row anchor points via DBSCAN clustering followed by RANSAC fitting to generate navigation lines, with reported angular and lateral errors minimized on the central crop row~\cite{sensors2025deeplabv3}. While functional in dense, uniform vegetation, this approach depends on continuous plant growth and degrades when seedlings are missing or canopy density is irregular. In addition, RANSAC is non-deterministic: random hypothesis generation introduces frame-to-frame angular jitter between consecutive video frames, requiring downstream temporal filtering to stabilize steering commands. The stochastic nature of RANSAC stems from its iterative random sampling strategy: each iteration selects a minimal subset of points (typically two for line fitting), computes a candidate model, and counts inliers within a residual threshold. To achieve $99\%$ confidence of finding an all-inlier sample for a two-point linear model when outliers constitute $50\%$ of the data, the algorithm requires $16$ iterations, rising to $35$ iterations when the outlier ratio reaches $65\%$; this sample count grows exponentially as outlier contamination increases~\cite{fischler1981ransac}. On embedded platforms operating at $20$~Hz with strict latency budgets, the computational overhead and output variance of RANSAC present practical limitations for closed-loop control.

    \item \textbf{Global contour tracking with line fitting:} An alternative approach applies edge operators (such as Canny filtering) or boundary tracing algorithms (\texttt{findContours}) over the full segmented corridor mask, separating the resulting pixel cloud into lateral boundaries before fitting lines with RANSAC. On dense and noisy boundary masks, this introduces variable computational latency and retains the stochastic jitter inherent to random sampling, confirming that raw RANSAC output on monocular imagery lacks the determinism required for smooth steering commands.

    \item \textbf{Horizontal strip decomposition with centroid reduction:} Methods operating on segmented drivable soil corridors (such as the multi-ROI scheme of Li et al.~\cite{li2022multiroi}) partition the image into horizontal bands and extract the bounding interval $[x_L, x_R]$ within each strip. However, these methods immediately collapse each interval into its midpoint:
    \begin{equation}
        x_C = \frac{x_L + x_R}{2}
    \end{equation}
    fitting curves (such as cubic B-splines or polynomials) directly to the sequence of centroids $x_C$. This premature reduction discards the distinction between the two physical edges. Under asymmetric disturbances---such as protruding leaves, localized weeds, or shadows extending along one side of the corridor---the midpoint shifts toward the intrusion, which introduces systematic bias in both lateral offset and heading angle.

    \item \textbf{Cascaded deep learning for line extraction:} Recent studies investigate two-stage deep architectures, where an initial segmentation model is followed by a secondary network dedicated to semantic line detection. While capable of learning line features directly from data, running two cascaded neural networks on embedded platforms exceeds the memory and latency limits required for closed-loop robotic control ($\ge 20$~Hz).

    \item \textbf{Single-edge elevation fitting:} M-estimators have also been applied to crop boundaries, such as Huber-loss Iteratively Reweighted Least Squares (IRLS) on stereo elevation scanlines~\cite{kneip2020cropedge}. However, these formulations track a single cutting boundary along dense harvesting fronts, rather than parallel crop boundaries defining an inter-row corridor from monocular RGB imagery.
\end{enumerate}

To address the trade-off between computational determinism and disturbance decoupling without cascading neural networks, this thesis investigates a decoupled dual-edge linear formulation:
\begin{equation}
    x_L(y) = m_L y + q_L, \qquad x_R(y) = m_R y + q_R
\end{equation}
Both boundaries are estimated by Iteratively Reweighted Least Squares (IRLS) with the weight rule
\begin{equation}
    w_i = \frac{1}{1 + \left(\frac{r_i}{\sigma}\right)^2}
\end{equation}
applied for a fixed number of four iterations to at most one boundary sample per horizontal strip ($N = 24$ by default); Chapter~\ref{cap:geometric_reference} details how \texttt{numpy.polyfit} applies these weights. A perturbation on one boundary changes only the weights and fit of that boundary. The centerline reference $x_C(y) = (x_L(y) + x_R(y))/2$ is computed after both boundary fits.

\subsection{Comparison of geometric representations}
\label{subsec:geometric_comparison}

Table~\ref{tab:confronto_riferimenti_geometrici} summarizes the input
representation and geometric output of the methods discussed above.

\begin{table}[htbp]
    \centering
    \footnotesize
    \setlength{\tabcolsep}{3pt}
    \renewcommand{\arraystretch}{1.15}
    \caption{Perceptual inputs and geometric references in selected
    agricultural navigation studies.}
    \label{tab:confronto_riferimenti_geometrici}
    \begin{tabularx}{\textwidth}{@{}>{\raggedright\arraybackslash}p{0.16\textwidth}
        >{\raggedright\arraybackslash}X
        >{\raggedright\arraybackslash}X
        >{\raggedright\arraybackslash}p{0.22\textwidth}@{}}
        \toprule
        Study & Perceptual input & Geometric operation & Output \\
        \midrule
        Chen et al.\ \cite{chen2021navigation}
        & Plant--soil segmentation and navigation points
        & Restricted Hough search around a predicted point
        & Image-space path line \\
        \addlinespace
        Kneip et al.\ \cite{kneip2020cropedge}
        & Stereo elevation map
        & Scanline transitions and IRLS linear fit
        & Ground-plane crop edge \\
        \addlinespace
        Mendez et al.\ \cite{mendez2023vineyards}
        & Thresholded background region
        & Column counts and mean of the five largest
        & Horizontal row-center position \\
        \addlinespace
        Navone et al.\ \cite{navone2025gpsfree}
        & Vegetation mask with depth
        & Smoothed column histogram; inverse-depth weighting in SegMinD
        & Histogram-minimum position \\
        \addlinespace
        Li et al.\ \cite{li2022multiroi}
        & Predicted inter-row-path mask
        & Adaptive horizontal regions and cubic B-spline
        & Curved image-space path \\
        \addlinespace
        De Silva et al.\ \cite{silva2023rowdetection}
        & Predicted crop-row mask
        & Upper anchor and triangular line search
        & Central crop-row line \\
        \midrule
        \addlinespace
        This thesis (Proposed DELE)
        & Predicted soil-corridor / mulch mask
        & Discrete horizontal strips, separate IRLS fits of the two boundaries
        & Straight image-space centerline $x_C(y)$ and $(e_y, e_\theta)$ \\
        \bottomrule
    \end{tabularx}
\end{table}

The methods are not directly ranked by this comparison: their inputs
and outputs differ. Li et al. provide the closest structural precedent
for strip-based point selection, while TSM illustrates a line search
for a crop-row representation. Neither method uses the same
soil-corridor mask and straight-line model as the pipeline developed
in this thesis.

\section{System integration and experimental validation}
\label{sec:system_integration}

\subsection{Perception and control integration}
\label{subsec:perception_control}

Navigation systems use the extracted reference to generate motion
commands. Navone et al. convert the displacement of their histogram
reference into linear and angular velocity commands and smooth them
with an Exponential Moving Average
\cite[Section~3.4]{navone2025gpsfree}. Mendez et al. use proportional
control based on row-center displacement and the distribution of
background pixels \cite[Section~3.5]{mendez2023vineyards}. De Silva et
al. combine the central row's angular and positional errors in a
proportional visual-servoing controller for simulation
\cite[Section~4.3]{silva2023rowdetection}.

\subsection{Scope of reported evaluations}
\label{subsec:evaluation_scope}

The studies evaluate different stages of navigation. Kneip et al.
assess crop-edge detection on recorded harvesting data
\cite[Section~8]{kneip2020cropedge}; Chen et al. report angular errors
and processing times for image-based path fitting
\cite[Sections~3.6, Tables~6--7]{chen2021navigation}. De Silva et al.
evaluate row detection on real images and visual servoing in simulation,
with preliminary real-world row following also reported
\cite[Sections~5.2--5.3 and~6, Table~5]{silva2023rowdetection}.
Mendez et al. report closed-loop vineyard trials, measuring centering
error from the distances to the two sides of the row
\cite[Sections~3.6 and~4]{mendez2023vineyards}. Navone et al. test in
simulation and crop fields, using LiDAR and attitude measurements to
assess displacement and heading \cite[Sections~4.3--4.4]{navone2025gpsfree}.

Segmentation IoU, image-space line error, field tracking error, and
processing time measure different parts of the system. Results from
these studies therefore provide context, not a common performance
benchmark.

\section{Positioning of this thesis}
\label{sec:thesis_positioning}

The system developed in this thesis addresses autonomous navigation in GNSS-denied crop rows using monocular RGB vision and embedded computation. Its technical positioning relative to the state of the art rests on three architectural decisions:

\begin{enumerate}
    \item \textbf{Perceptual representation (Soil corridor versus canopy):} Rather than detecting crop stalks or foliage masks, the segmentation stage targets the navigable soil corridor (and plastic mulch). This makes the navigation reference independent of crop growth stage, canopy height, and missing plants (fallanze), providing a continuous geometric guide directly at ground level.
    
    \item \textbf{Geometric extraction (Decoupled deterministic boundaries versus centroid reduction):} Rather than collapsing strip intervals into midpoints prior to fitting (as in Li et al.~\cite{li2022multiroi}) or employing stochastic RANSAC over dense edge contours, the proposed pipeline decouples the left and right physical corridor boundaries. Each boundary is fitted independently by iteratively reweighted least squares with a fixed number of four iterations. A disturbance on one side (such as a leaf overhang or a shadow) enters only the fit of that boundary; the centerline, as the mean of the two edge lines, receives half of that boundary's deviation.
    
    \item \textbf{Embedded deployment feasibility:} The architecture decouples offline automatic dataset generation from onboard execution. Large foundation models (SAM~3) are restricted to offline zero-shot concept annotation, while real-time inference on the NVIDIA Jetson Orin Nano is handled by a lightweight LR-ASPP network (using a MobileNetV3 backbone) paired with the deterministic IRLS geometric extractor. The entire pipeline executes in closed loop at 20~Hz without requiring secondary neural networks, stereo elevation mapping, or LiDAR sensors.
\end{enumerate}

The chapters that follow detail and evaluate each stage of this pipeline: Chapter~\ref{cap:architectures} formulates the deep learning architectures, Chapter~\ref{cap:segmentation} analyzes the annotation protocols and comparative model training, Chapter~\ref{cap:geometric_reference} formalizes the geometric extractors, Chapter~\ref{cap:platform} describes the Demetra robotic platform and onboard ROS integration, and Chapter~\ref{cap:experiments} reports offline benchmarks and field tests.
